\documentclass{jsarticle}
\usepackage[dvipdfmx, final]{graphicx}
\usepackage{amsmath}
\usepackage{amssymb}
\usepackage{chemfig}
\begin{document}
\title{Zeta function escimat with beta function of remake formula}
\author{Masaaki Yamaguchi}
\date{}
\maketitle
ベータ関数は、
$$ \beta(p,q) = {{\Gamma(p)\Gamma(q)} \over {\Gamma(p+q)}} $$
このゼータ関数版が、
$$ {\zeta(p)\cdot\zeta(q) \over {\zeta(p+q)}} = \bigoplus{(
i\hbar^{\nabla})}^{\oplus L} = e^{-x\log x}$$
$y=x\log x$がゼータ関数で、$z=e^{x\log x}$もゼータ関数で、$e^{-f}dV=
{1 \over {\log x}}$が素数になっていて、$\zeta(s)={{\beta(p,q)}\over 
{\log x}}$もゼータ関数で、$\text{素数}={\text{偶数}\over\text{偶数}}$
が$\sqrt{2}$と同じ仕組みになっていて、
$y=x^{\log x}$と$z=e^{x\log x}$はともにゼータ関数であり、
$e^{-f}dV={1 \over {\log x}}$は素数分布に関係している。
$\zeta(s)={{\beta(p,q)}\over {\log x}}$はゼータ関数の形をしている。
$\text{素数}={\text{偶数}\over\text{偶数}}$は、$\sqrt{2}$ 
と同様の構造をしている。
$\int e^{-f}dV=\int\Gamma^{'}(\gamma)dx_m={\Gamma \over {\log x}}$
と、サーストン・ペレルマン多様体である、整数を$\log x$の奇数で
割った結果の偶数が整数の補空間がゼータ関数の奇数性により、
　$$ \int e^{-f}dV = \int \Gamma^{'}(\gamma)dx_m = {\Gamma\over{\log x}} $$
これから、
$$ ||ds^2|| = 8\pi G\left({p \over c^3} + {V \over S}\right) $$
この3次元多様体のエントロピー値が、
$$ \nabla^{2} \Psi = 4\pi G(\rho + {p \over c^2}) $$
$$ \nabla^{2} \Psi = 4 \pi G(\rho + {p \over c^2}) + K{\zeta(2n+1) \over 
{\log x}} $$
$$ 8\pi G(\rho + {p \over c^2}) + K{\zeta(2n+1) \over {\log x}}
={R^2 \over {z^2}} $$
$$ ||ds^2|| = e^{R + {1 \over 2}g_{ij}\Lambda
\log{( R + {1 \over 2}g_{ij}\Lambda)}}
+ {K\zeta(2n+1) \over {\log x}}$$
と言えて、
$$ ||ds^2|| = e^{-2\pi T||\psi||}[\eta + \bar{h}]dx^{\mu}dx^{\nu} +
T^2 d^2 \psi $$
と、AdS5多様体となる。

指数関数における、相加相乗平均方程式が、
$$ \log Z = \log{{p^q\cdot q^p} \over {{(p+q)}^{(p+q)}}} $$
$$ p + q = {}^{p+q}\sqrt{p^q\cdot q^p} $$
$$ = {}^q\sqrt{p}\cdot {}^p\sqrt{q} $$
ゼータ関数におけるシャノンの公式を、非可換代数多様体に使うと、
$$ \pi(\chi,x) = [i\pi(\chi,x),f(x)] $$
$$ p=x,q=y \to x + y = e^{-x\log x}\cdot e^{-y\log y} $$
$$ = e^{-(x\log x + y \log y)} $$
$$ = \int e^{-x^2 - y^2}dxdy = \pi $$
という、ベータ関数のゼータ関数版が、ガンマ関数になり、
ガウス関数と帰結して、サーストン・ペレルマン多様体である、
ヒルベルト空間で、全ての数学モデルが成り立っている。


$$ {{e^{x\log x + y\log y}} \over {e^{x\log(x+y) + y\log(x+y)}}} $$
$$ = z $$
サラスの公式である、上の数式展開は、素粒子方程式である、
益川・小林理論の帰結でもある。

\end{document}


